\section{POMDP Model}
\label{POMDP-model}


A preliminary version of our POMDP model appeared at the SecArt'11
workshop \cite{SarBufHof11}. The reader may refer to that paper for a
more detailed example listing complete transition and observation
models for some actions, and exemplifying the evolution of belief
states when applying these actions. In what follows, we keep the
description brief in the interest of space.





\subsection{States}
\label{POMDP-model:states}


Several aspects of the problem---notably the network structure and the
firewall filtering rules---are known and static. POMDP variables
encoding these aspects can be compiled out in a pre-process, and are
not included in our model.

The states capture the status of each machine (controlled/reached/not
reached). For non-controlled machines, they also specify the software
configuration (operating system, servers, open ports, \dots). We
specify the vulnerable programs, as well as programs that can provide
information about these (e.g., the protection mechanism ``DEP'' in our
running example is relevant to both exploits). The states also
indicate whether a given machine or program has crashed. 

Finally, we introduce one special \emph{terminal} state into the POMDP
model (of the entire network, not of individual machines). That state
corresponds to giving up the attack, when for every available action
(if any) the potential benefit is not worth the action's cost.


\vspace{-0.05cm}
\begin{example}\label{exp:running-states}
The states describe the attacked machine $m$. For simplicity, we
assume that the exploits here do not risk crashing the machine (see
also next sub-section). Apart from the terminal state and the state
representing that $m$ is controlled, the states specify which programs
(``SA'' or ``CAU'') are present, whether they are vulnerable, and
whether ``DEP'' is enabled. Each application is listening on a
different port, so a port is open iff the respective application is
present, and we do not need to model ports separately. Thus we have a
total of $20$ states:\\ {\scriptsize
\begin{minipage}[t]{.24\linewidth}
\begin{verbatim} 
1 terminal
2 m_controlled
\end{verbatim}
\end{minipage}
% \hfill
\begin{minipage}[t]{.34\linewidth}
\begin{verbatim}
 3 m_none
 4 m_CAU
 5 m_CAU_Vul
 6 m_SA
 7 m_SA_CAU
 8 m_SA_CAU_Vul
 9 m_SA_Vul
10 m_SA_Vul_CAU
11 m_SA_Vul_CAU_Vul
\end{verbatim}
\end{minipage}
%  \hfill
\begin{minipage}[t]{.30\linewidth}
\begin{verbatim}
12 m_DEP_none
13 m_DEP_CAU
14 m_DEP_CAU_Vul
15 m_DEP_SA
16 m_DEP_SA_CAU
17 m_DEP_SA_CAU_Vul
18 m_DEP_SA_Vul
19 m_DEP_SA_Vul_CAU
20 m_DEP_SA_Vul_CAU_Vul
\end{verbatim}
\end{minipage}
}
\end{example}
%\vspace{-0.05cm}



In short, the states for each machine $m$ essentially are tuples of
status values for each relevant program. Global system states then are
tuples of these machines-states, with one entry for each $m \in
M$. The state space enumerates these tuples. In other words, the state
space is factored in a natural way, by programs and machines. An
obvious option is, thus, to model and solve the problem using factored
POMDPs (e.g., \cite{HanFen00}). We did not try this yet; our POMDP
model generator internally enumerates the states, and feeds the ground
model to SARSOP.\footnote{Note that this approach enables certain
  non-trivial optimizations: some of the states in
  Example~\ref{exp:running-states} could be merged. If DEP is enabled,
  then it does not matter whether or not CAU/SA are vulnerable. For
  brevity, we do not discuss this in detail here.}


The factored nature of our problem also implies that the state space
is huge. In a realistic setting, the set $C$ of possible configuration
tuples for each machine $m \in M$ is very large, yielding an enormous
state space $|\cS|=O(|C|^{|M|})$. In practice, we will run POMDPs only
on single machines, i.e., $|M|=1$.




\subsection{Actions}
\label{POMDP-model:actions}

To reach the terminal state, we need a \emph{terminate} action
indicating that one gives up on the attack. 

There are two main types of actions, \emph{scans} and \emph{exploits},
which both have to be targeted at reachable machines.
%
Scans can be OS detection actions or port scans. In most cases, they
have no effect on the state of the target machine. Their purpose is to
gain knowledge about a machine's configuration, by an observation that
typically allows to prune some states from the belief (e.g., observing
that the OS must be some Windows XP version). %
Exploits make use of a vulnerability---if present---to gain control
over a machine. The outcome of the exploit is observed by the
attacker, so a failed exploit may, like a scan, yield information
about the configuration (e.g., that a protection mechanism is likely
to be running). For a minority of exploits, a failed attempt crashes
the machine. 

For all actions, the outcome is deterministic: which observation is
returned, and whether an exploit succeeds/fails/crashes, is uniquely
determined by the target machine's configuration.

%\vspace{-0.05cm}
\begin{example}\label{exp:running-actions}
In our example, there are five possible actions:\\ {\scriptsize
\begin{minipage}[t]{.4\linewidth}
\begin{verbatim}
m_exploit_SA
m_exploit_CAU
m_scan_port_2967
m_scan_port_6668
terminate
\end{verbatim}
\end{minipage}
}

\smallskip
\noindent
The POMDP model specifies, for each state in
Example~\ref{exp:running-states}, the outcome of each action. For
example, {\scriptsize \verb+m_exploit_SA+} succeeds if and only if SA
is present and vulnerable, and DEP is disabled. Hence, when applied to
either of the states 9, 10, or 11, {\scriptsize \verb+m_exploit_SA+}
results in state 2, and returns the observation {\scriptsize
  \verb+succeeded+}. Applied to any other state, {\scriptsize
  \verb+m_exploit_SA+} leaves the state unchanged, and the observation
is {\scriptsize \verb+failed+}.
\end{example}
%\vspace{-0.05cm}

The outcomes of actions also depend on what firewall (if any) stands
between the pentester and the target. If the firewall filters out the
relevant port, then the action is unusable: its transition model
leaves the state unchanged, and no observation is returned. For
example, if a firewall $F$ filters out port 2967, then {\scriptsize
  \verb+m_scan_port_2967+} and {\scriptsize \verb+m_exploit_SA+} are
unusable through $F$, but can be employed as soon as a machine behind
$F$ is under control.




\subsection{Rewards}
\label{POMDP-model:rewards}

No reward is obtained when using the \emph{terminate} action or when
in the terminal state. 

The instant reward of any scan/exploit action depends on the
transition it induces in the present state. Our simple model is to
additively decompose the instant reward $\rewardfn(s,a,s')$ into
$\rewardfn(s,a,s') = r_e(s,a,s') + r_t(a) + r_d(a)$.  Here, (i)
$r_e(s,a,s')$ is the value of the attacked machine in case the
transition $(s,a,s')$ corresponds to a successful exploit, and is $0$
for all other transitions; (ii) $r_t(a)$ is a cost that depends on the
action's duration; %
and (iii) $r_d(a)$ is a cost that reflects the risk of detection when
using this action. (iii) is orthogonal to the risk of crashing a
program/machine, which as described we model as a possible outcome of
exploits. Note that (ii) and (iii) may be correlated; however, there
is no 1-to-1 correspondence between the duration and detection risk of
an exploit, so it makes sense to be able to distinguish these
two. Finally, note that (i) results in summing up rewards for
successful exploits on different machines. That is not a limiting
assumption: one can reward breaking into \emph{[$m_1$ OR $m_2$]} by
introducing a new virtual machine, accessible at no cost from each of
$m_1$ and $m_2$.


\vspace{-0.05cm}
\begin{example}\label{exp:running-rewards}
In our example, we set $r_e=100$ in case of success, $0$ otherwise;
$r_t=-10$ for all actions; and $r_d=0$ (no risk of detection). We will
see below what effect these settings have on an optimal policy.
\end{example}
\vspace{-0.05cm}

Since all actions are deterministic, there is no point in repeating
them on the same target through the same firewall---this will not
produce new effects or bring any new information. In particular,
positive rewards cannot be received multiple times. Thus cyclic
behaviors incur infinite negative costs. This implies that the
expected reward of an optimal policy is finite even without
discounting.\footnote{In fact, the problem falls into the class of
  \emph{Stochastic Shortest Path Problems} \cite{BerTsi96}.}




\subsection{Designing the Initial Belief}
\label{POMDP-model:initial-belief}


Penetration testing is done at regular time intervals. The initial
belief---our knowledge of the network when we start the
pentesting---depends on (a) what was known at the end of the previous
pentest, and on (b) what may have changed since then. We assume for
simplicity that knowledge (a) is perfect, i.e., each machine $m$ at
time $0$ (the last pentest) is assigned one concrete configuration
$I(m)$. We then compute the initial belief as a function $b_0(I,T)$
where $T$ is the number of days elapsed since the last pentest. The
uncertainty in this belief arises from not knowing which software
updates were applied. We assume that the updates are made
independently on each machine (simplifying, but reasonable given that
updates are controlled by individual users).

A simple model of updates \cite{SarBufHof11} encodes the uncertain
evolution of each program independently, in terms of a Markov
chain. The states in each chain correspond to the different versions
of the program, and the transitions model the possible program updates
(with estimated probabilities that these updates will be made). The
initial belief then is the distribution resulting from this chain
after $T$ steps.

\vspace{-0.05cm}
\begin{example}\label{exp:generating-initial-belief}

  \begin{figure}[t]

    \begin{minipage}{0.3\linewidth}
      \centerline{\includegraphics[width=0.95\linewidth]{mc-DEP.pdf}}
    \end{minipage}
    \hfill
    \begin{minipage}{0.3\linewidth}
      \centerline{\includegraphics[width=0.95\linewidth]{mc-SA.pdf}}
    \end{minipage}
%
%    \medskip
\hfill
%    \centerline{
      \begin{minipage}{0.3\linewidth}
        \centerline{\includegraphics[width=0.95\linewidth]{mc-CAU.pdf}}
      \end{minipage}
%    }
\vspace{-0.2cm}
    \caption{The three independent Markov chains used to model the
      update mechanism in our example network.}
    \label{fig:MarkovChains}
\vspace{-0.6cm}
  \end{figure}

  In our running example, the three components in the single machine
  are DEP, CAU and SA. They are updated via three independent Markov
  chains, each with two states, as illustrated in
  Figure~\ref{fig:MarkovChains}. The probabilities indicate how likely
  the machine is to transition from one state to another during one
  day. Say we set $T=30$, and run the Markov chains on the
  configuration $I$ in which $m$ has DEP disabled, and both SA and CAU
  are vulnerable to the attacker's exploit. In the resulting initial
  belief $b_0(I,T)$, DEP is likely to be enabled; the weight of states
  12--20 in Example~\ref{exp:running-states} is high in $b_0$
  ($>70\%$).
\end{example}
\vspace{-0.05cm}

Here, we use this simple model as the basic building block in a method
taking into account that version $x$ of program A may need version $y$
or $z$ of program B. We assume that programs are organized in a
hierarchical manner, the operating system being at the root of a
directed acyclic graph, and a program having as its parents the
programs it directly depends on. This yields a Dynamic Bayesian
Network, where each conditional probability distribution is derived
from a Markov chain $Pr(X_t=x'|X_{t-1}=x)$ %
filtered by a compatibility function $\delta(X=x,parent_1(X)=x_1,
\dots, parent_k(X)=x_k)$, that returns $1$ iff the value of $X$ is
compatible with the parent versions, $0$ otherwise. This model of
updates is reasonable, but of course still not realistic; future work
needs to investigate such models in detail.



We now illustrate how reasoning with the probabilities of the initial
belief results in the desired intelligent behavior.

\vspace{-0.05cm}
\begin{example}\label{exp:running-initial-belief-policy}
Say we compute the initial belief $b_0(I,T)$ as in
Example~\ref{exp:generating-initial-belief}. Since the weight of
states 12--20 is high in $b_0$, if {\scriptsize \verb+m_exploit_SA+}
fails, then the success probability of {\scriptsize
  \verb+m_exploit_CAU+} is reduced to the point of not being worth the
effort anymore, and the attacker (the optimal policy) gives up, i.e.,
would try a different attack not prevented by DEP. Namely, consider
$Pr($CAU$^+|$2967$^+)$, i.e., the probability of {\scriptsize
  \verb+m_exploit_CAU+} succeeding, after observing that port 2967 is
open. This corresponds to the weight of (A) states 8 and 11 in
Example~\ref{exp:running-states}, within the states (B) 6--11 plus
15--20. That weight (A/B) is about $20\%$. Thus the expected value of
{\scriptsize \verb+m_exploit_CAU+} in this situation is about
$100*0.2$ [success reward] $- 10$ [action cost] $= 10$,
cf.\ Example~\ref{exp:running-rewards}, so the action is
worthwhile. By contrast, say that {\scriptsize \verb+m_exploit_SA+}
has been tried and failed. Then (A) is reduced to state 8 only, while
(B) still contains (in particular) all the DEP states 15--20. The
latter states have a lot of weight, and thus
$Pr($CAU$^+|$2967$^+,$SA$^-)$ is only about $5\%$. Given this, the
expected value of {\scriptsize \verb+m_exploit_CAU+} is negative, and
it is better to apply {\scriptsize \verb+terminate+} instead.
\end{example}
\vspace{-0.05cm}


